% !TEX root =  Lecture11.tex


%%%%%%%%%%%%%%%%%%%%%%%%%%%%%%%%%%%%
% Recap/Agenda
%%%%%%%%%%%%%%%%%%%%%%%%%%%%%%%%%%%%
\begin{frame}
\frametitle{Module 2: Statistical Inference}
    \begin{itemize}
        \item \hl{Previously:} confidence intervals by \emph{bootstrapping} --- resample the data, then read the interval off the simulated distribution.
        \item \hl{This time:} the same intervals and p-values, from \emph{formulas}.
          \begin{itemize}
              \item exact calculations: the binomial and the normal model;
              \item the central limit theorem --- and why our simulated distributions kept coming out bell-shaped;
              \item p-values and confidence intervals with $z$ and $t$;
              \item two-sample and paired data.
          \end{itemize}
        \item \hl{Reading:} IMS Chapter 13.
        \item \hl{Homework for today:}
        \item \hl{Homework for next class:}
    \end{itemize}
\end{frame}


%%%%%%%%%%%%%%%%%%%%%%%%%%%%%%%%%%%%
% Learning objectives:
%%%%%%%%%%%%%%%%%%%%%%%%%%%%%%%%%%%%
\begin{frame}
    \frametitle{Lecture Objectives}
    \goals{%
    \begin{itemize}
    \item compute an exact binomial probability and a $z$-score;
    \item state the central limit theorem and use it instead of simulating;
    \item build a CI and p-value from a formula, and why they match simulation;
    \item handle two-sample and paired problems, and recognize paired data;
    \item choose between a formula and a simulation, and say what each costs.
    \end{itemize}}
    \objectives{%
    \begin{itemize}
    \item \textbf{(M2, LO1)} Approaches to Inference
    \item \textbf{(M2, LO2)} Simulation-based Inference
    \end{itemize}}
\end{frame}
